% Laget automatisk av elevark/lag_elevark.py fra sim/bremselengde.html. Endringer her blir overskrevet.
\documentclass[11pt]{article}
\usepackage{elevark}
\begin{document}
\elevarkhode{Fysikk 1 · Rettlinjet bevegelse}{Reaksjons- og bremselengde}{https://fysikklab.no/sim/bremselengde.html}{Simuleringen}

Et barn løper ut i veien foran en bil. Først kjører bilen videre med samme fart mens sjåføren reagerer. Så bremser den med konstant akselerasjon. Stopper bilen i tide? Og hva skjer hvis den kjørte bare litt fortere?

\section*{Oppgaver}

Start med 50~km/h, reaksjonstid 1,0~s, tørr asfalt og 30~m til hinderet.

\begin{oppgaveliste}
\oppgave \fase{test}{Regn}Regn ut reaksjonslengden, bremselengden og stopplengden. Stopper bilen før hinderet? Sjekk med simuleringen.
\svar
\oppgave \fase{gjett}{Gjett}Bilen kjører i stedet 60~km/h. Stopper den før hinderet? Hvis ikke, hvor fort tror du den kjører når den treffer?
\svar
\oppgave \fase{test}{Regn}Regn ut farten bilen treffer hinderet med i 60~km/h. Hvordan kan 10~km/h ekstra gi en så stor forskjell?
\svar
\oppgave \fase{forklar}{Forklar}Hva skjer med reaksjonslengden og bremselengden når farten dobles, fra 50 til 100~km/h? Forklar ut fra formlene.
\svar
\oppgave \fase{test}{Undersøk}Prøv våt asfalt, snø og is med 50~km/h. Hvor fort kan du kjøre på is hvis bilen skal stoppe før 30~m?
\svar
\oppgave \fase{test}{Undersøk}En sjåfør som ser på mobilen, kan få reaksjonstid på 2~s eller mer. Hva skjer med stopplengden ved 50~km/h?
\svar
\oppgave \fase{utfordring}{Utfordring}Lag en formel for den største farten bilen kan ha for å stoppe før hinderet, uttrykt med $D$, reaksjonstiden $t_{\mathrm{r}}$ og $a$. Bruk den for tørr asfalt, 1,0~s og 30~m.
\svar
\end{oppgaveliste}

\end{document}
